\documentclass[11pt,a4paper]{article}
\usepackage[T1]{fontenc}
\usepackage[utf8]{inputenc}
\usepackage{amsmath,amssymb,xcolor}
\usepackage[margin=25mm]{geometry}
\usepackage[colorlinks=true,urlcolor=blue]{hyperref}
\definecolor{nexorablue}{HTML}{2459B5}
\setlength{\parindent}{0pt}
\setlength{\parskip}{7pt}
\setlength{\emergencystretch}{2em}
\renewcommand{\familydefault}{\sfdefault}
\begin{document}

{\small\bfseries\color{nexorablue} NEXORA CPGE}\par
{\small MPSI / NEX-RGT-03 / Corrigé}\par
{\LARGE\bfseries Charge depuis une tension non nulle}\par
\vspace{0.25em}\hrule\vspace{0.6em}

\section*{Équation et état initial}
On a $\tau=RC=0{,}47\,\mathrm s$, puis $u_C(t)=6-4e^{-t/\tau}$ en volts. À $0^+$, le courant vaut $(6-2)/R=0{,}40\,\mathrm{mA}$.

\section*{Franchissement}
La tension atteint $5\,\mathrm V$ quand $e^{-t/\tau}=1/4$, soit $t=0{,}47\ln4\simeq0{,}65\,\mathrm s$. La tension croît et le seuil est inférieur à sa limite, donc l’instant est unique.


\end{document}
